% Part: history
% Chapter: biographies
% Section: ernst-zermelo

\documentclass[../../../include/open-logic-section]{subfiles}

\begin{document}

\olfileid[tr]{his}{bio}{zer}

\olsection{Ernst Zermelo}

\olphoto{zermelo-ernst}{Ernst Zermelo}

Ernst Zermelo, 27 Temmuz 1871'de Almanya'nın Berlin kentinde doğdu. Beş kız
kardeşi vardı; ancak ailesi sağlık sorunları yaşadı ve kardeşlerin yalnızca
üçü yetişkinliğe erişebildi. Anne babası da Zermelo küçükken öldü; o ve
kardeşleri, Zermelo on yedi yaşındayken öksüz kaldılar. Zermelo sanata,
özellikle şiire büyük ilgi duyuyordu. Keskin, nükteli ve eleştirel olmasıyla
tanınırdı. En ünlü matematik başarıları arasında seçim aksiyomunu ortaya
atması (1904) ve küme teorisini aksiyomlaştırması (1908) yer alır.

Zermelo'nun üniversitedeki ilgi alanları çeşitlilik gösteriyordu. Fizik,
matematik ve felsefe dersleri aldı. Hermann Schwarz'ın danışmanlığında
\emph{Varyasyonlar Hesabında İncelemeler} başlıklı tezini 1894'te Berlin
Üniversitesi'nde tamamladı. 1897'de G\"{o}ttingen Üniversitesi'nde daha fazla
öğrenim görmeye karar verdi; burada David Hilbert'in temellere ilişkin
çalışmalarından büyük ölçüde etkilendi. 1899'da profesörlüğe hak kazandı;
ancak---belki de tuhaf davranışları ve ``sinirli telaşı'' nedeniyle---on bir
yıl boyunca profesörlük alamadı.

Zermelo nihayet 1910'da Zürih Üniversitesi'nde ücretli bir profesörlük aldı;
ancak verem nedeniyle 1916'da emekliye ayrılmak zorunda kaldı. İyileşmesinin
ardından 1921'de Freiburg Üniversitesi'nde fahri profesörlük verildi. Bu
dönemde matematiğin temelleri üzerinde çalıştı. Thoralf Skolem ve Kurt
G\"{o}del'in çalışmalarına öfkelendi ve makalelerinde onların yaklaşımlarını
açıkça eleştirdi. Gözden düşmesi ve Hitler'in Almanya'da iktidara yükselişine
karşı çıkması nedeniyle 1935'te Freiburg'daki görevinden çıkarıldı.
 
Zermelo'nun yaşamının son yıllarına yalnızlık damgasını vurdu. 1935'te
görevden çıkarıldıktan sonra matematiği bıraktı. Kırsal bölgeye taşındı ve
mütevazı biçimde yaşadı. 1944'te evlendi; görme yetisini yitirmekte olduğu
için eşine bütünüyle bağımlı hâle geldi. Zermelo 1951'e gelindiğinde görme
yetisini tamamen kaybetmişti. 21 Mayıs 1953'te Almanya'nın G\"{u}nterstal
kentinde öldü.

\begin{reading}
Zermelo'nun tam yaşam öyküsü için \citet{Ebbinghaus2015} kaynağına bakın.
Zermelo'nun çığır açıcı 1904 ve 1908 tarihli makaleleri özgün Almancalarıyla
okunabilir \citep{Zermelo1904,Zermelo1908}. Zermelo'nun fizik yazıları dâhil
toplu eserleri İngilizce çeviri olarak \citep{Ebbinghaus2010,Ebbinghaus2013}
içinde mevcuttur.
\end{reading}

\end{document}
